\documentclass[11pt]{article}
\usepackage[letterpaper,margin=0.83in,headheight=16pt]{geometry}
\usepackage[T1]{fontenc}
\usepackage[utf8]{inputenc}
\usepackage{lmodern}
\usepackage{microtype}
\usepackage{amsmath,amssymb}
\usepackage{xcolor}
\usepackage{hyperref}
\usepackage{enumitem}
\usepackage{booktabs,longtable,tabularx,array}
\usepackage{fancyhdr}
\usepackage{titlesec}
\usepackage{listings}
\usepackage{tikz}
\usepackage{float}
\definecolor{suiteblue}{HTML}{17365D}
\definecolor{suitelight}{HTML}{EAF1F8}
\definecolor{suitegray}{HTML}{556070}
\hypersetup{colorlinks=true,linkcolor=suiteblue,urlcolor=suiteblue,breaklinks=true}
\titleformat{\section}{\Large\bfseries\color{suiteblue}}{\thesection}{0.55em}{}
\titleformat{\subsection}{\large\bfseries\color{suiteblue}}{\thesubsection}{0.55em}{}
\setlist[itemize]{itemsep=0.25em,topsep=0.35em,leftmargin=1.4em}
\setlist[enumerate]{itemsep=0.25em,topsep=0.35em,leftmargin=1.6em}
\setlength{\parindent}{0pt}
\setlength{\parskip}{0.47em}
\emergencystretch=2em
\lstset{basicstyle=\ttfamily\footnotesize,breaklines=true,frame=single,columns=fullflexible,showstringspaces=false}
\pagestyle{fancy}\fancyhf{}
\lhead{Game Design \& Engineering}\rhead{Workbook Template}\cfoot{\thepage}
\begin{document}
{\huge\bfseries\color{suiteblue} Probability And Balancing Workbook -- Template}\par
\vspace{0.15in}
{\normalsize\color{suitegray} Fictional example annotations plus reusable analytical prompts}\par
\vspace{0.15in}\hrule\vspace{0.15in}

\section{Record control}
Project: [Project]. Revision: [Revision]. Practitioner: [Name]. Evidence repository: [URI].


\section{Evidence discipline}
For each exercise, distinguish an assumption, result established by computation, result established by player observation, and an interpretation. The example entries are instructional and are not observations from a real game.


\section{Expected damage}
Exercise prompt: Weapon A deals three with probability 0.85. Weapon B deals six with probability 0.45. Compute mean and variance; identify what these values omit.

Expected deliverable: Hand calculation and simulation check.


\subsection{Hypothesis}
Illustrative example: A small design change affects a specified, observable player action.

Project-specific response: [Document the evidence and rationale.]


\subsection{Assumptions}
Illustrative example: The simulated population and state constraints are explicitly defined.

Project-specific response: [Document the evidence and rationale.]


\subsection{Model or procedure}
Illustrative example: Use a repeatable implementation and record input parameters.

Project-specific response: [Document the evidence and rationale.]


\subsection{Results and limitations}
Illustrative example: Quantitative result is supported by a run log; enjoyment remains unverified.

Project-specific response: [Document the evidence and rationale.]


\subsection{Decision and next test}
Illustrative example: Keep, revise, or retire the feature after a new independent study.

Project-specific response: [Document the evidence and rationale.]


\section{Dominance threshold}
Exercise prompt: Solve for accuracy p such that a five-damage weapon matches expected value of a three-damage weapon at accuracy 0.8.

Expected deliverable: Algebra, threshold, and sensitivity table.


\subsection{Hypothesis}
Illustrative example: A small design change affects a specified, observable player action.

Project-specific response: [Document the evidence and rationale.]


\subsection{Assumptions}
Illustrative example: The simulated population and state constraints are explicitly defined.

Project-specific response: [Document the evidence and rationale.]


\subsection{Model or procedure}
Illustrative example: Use a repeatable implementation and record input parameters.

Project-specific response: [Document the evidence and rationale.]


\subsection{Results and limitations}
Illustrative example: Quantitative result is supported by a run log; enjoyment remains unverified.

Project-specific response: [Document the evidence and rationale.]


\subsection{Decision and next test}
Illustrative example: Keep, revise, or retire the feature after a new independent study.

Project-specific response: [Document the evidence and rationale.]


\section{Sample uncertainty}
Exercise prompt: Estimate success probability with 100, 1000, and 10000 independently simulated encounters. Compare standard errors under the binomial assumption.

Expected deliverable: Seeded experiments, uncertainty summary, interpretation.


\subsection{Hypothesis}
Illustrative example: A small design change affects a specified, observable player action.

Project-specific response: [Document the evidence and rationale.]


\subsection{Assumptions}
Illustrative example: The simulated population and state constraints are explicitly defined.

Project-specific response: [Document the evidence and rationale.]


\subsection{Model or procedure}
Illustrative example: Use a repeatable implementation and record input parameters.

Project-specific response: [Document the evidence and rationale.]


\subsection{Results and limitations}
Illustrative example: Quantitative result is supported by a run log; enjoyment remains unverified.

Project-specific response: [Document the evidence and rationale.]


\subsection{Decision and next test}
Illustrative example: Keep, revise, or retire the feature after a new independent study.

Project-specific response: [Document the evidence and rationale.]


\section{Challenge and fairness}
Exercise prompt: Simulate two skill profiles and compare how the same rules alter time-to-completion and variance.

Expected deliverable: Per-profile comparison; non-metric fairness caveat.


\subsection{Hypothesis}
Illustrative example: A small design change affects a specified, observable player action.

Project-specific response: [Document the evidence and rationale.]


\subsection{Assumptions}
Illustrative example: The simulated population and state constraints are explicitly defined.

Project-specific response: [Document the evidence and rationale.]


\subsection{Model or procedure}
Illustrative example: Use a repeatable implementation and record input parameters.

Project-specific response: [Document the evidence and rationale.]


\subsection{Results and limitations}
Illustrative example: Quantitative result is supported by a run log; enjoyment remains unverified.

Project-specific response: [Document the evidence and rationale.]


\subsection{Decision and next test}
Illustrative example: Keep, revise, or retire the feature after a new independent study.

Project-specific response: [Document the evidence and rationale.]


\section{Economy breakage}
Exercise prompt: Construct one exploit in an uncapped source/sink model and fix it with a transaction invariant.

Expected deliverable: Exploit reproduction, corrected rules, regression test.


\subsection{Hypothesis}
Illustrative example: A small design change affects a specified, observable player action.

Project-specific response: [Document the evidence and rationale.]


\subsection{Assumptions}
Illustrative example: The simulated population and state constraints are explicitly defined.

Project-specific response: [Document the evidence and rationale.]


\subsection{Model or procedure}
Illustrative example: Use a repeatable implementation and record input parameters.

Project-specific response: [Document the evidence and rationale.]


\subsection{Results and limitations}
Illustrative example: Quantitative result is supported by a run log; enjoyment remains unverified.

Project-specific response: [Document the evidence and rationale.]


\subsection{Decision and next test}
Illustrative example: Keep, revise, or retire the feature after a new independent study.

Project-specific response: [Document the evidence and rationale.]


\section{Review decision}
Reviewer: [Name]. Findings: [Finding IDs]. Status: [Draft, revision needed, or accepted].


\end{document}
